
%% bare_conf.tex
%% V1.3
%% 2007/01/11
%% by Michael Shell
%% See:
%% http://www.michaelshell.org/
%% for current contact information.
%%
%% This is a skeleton file demonstrating the use of IEEEtran.cls
%% (requires IEEEtran.cls version 1.7 or later) with an IEEE conference paper.
%%
%% Support sites:
%% http://www.michaelshell.org/tex/ieeetran/
%% http://www.ctan.org/tex-archive/macros/latex/contrib/IEEEtran/
%% and
%% http://www.ieee.org/

%%*************************************************************************
%% Legal Notice:
%% This code is offered as-is without any warranty either expressed or
%% implied; without even the implied warranty of MERCHANTABILITY or
%% FITNESS FOR A PARTICULAR PURPOSE! 
%% User assumes all risk.
%% In no event shall IEEE or any contributor to this code be liable for
%% any damages or losses, including, but not limited to, incidental,
%% consequential, or any other damages, resulting from the use or misuse
%% of any information contained here.
%%
%% All comments are the opinions of their respective authors and are not
%% necessarily endorsed by the IEEE.
%%
%% This work is distributed under the LaTeX Project Public License (LPPL)
%% ( http://www.latex-project.org/ ) version 1.3, and may be freely used,
%% distributed and modified. A copy of the LPPL, version 1.3, is included
%% in the base LaTeX documentation of all distributions of LaTeX released
%% 2003/12/01 or later.
%% Retain all contribution notices and credits.
%% ** Modified files should be clearly indicated as such, including  **
%% ** renaming them and changing author support contact information. **
%%
%% File list of work: IEEEtran.cls, IEEEtran_HOWTO.pdf, bare_adv.tex,
%%                    bare_conf.tex, bare_jrnl.tex, bare_jrnl_compsoc.tex
%%*************************************************************************

% *** Authors should verify (and, if needed, correct) their LaTeX system  ***
% *** with the testflow diagnostic prior to trusting their LaTeX platform ***
% *** with production work. IEEE's font choices can trigger bugs that do  ***
% *** not appear when using other class files.                            ***
% The testflow support page is at:
% http://www.michaelshell.org/tex/testflow/



% Note that the a4paper option is mainly intended so that authors in
% countries using A4 can easily print to A4 and see how their papers will
% look in print - the typesetting of the document will not typically be
% affected with changes in paper size (but the bottom and side margins will).
% Use the testflow package mentioned above to verify correct handling of
% both paper sizes by the user's LaTeX system.
%
% Also note that the "draftcls" or "draftclsnofoot", not "draft", option
% should be used if it is desired that the figures are to be displayed in
% draft mode.
%
\documentclass[conference]{IEEEtran}
% Add the compsoc option for Computer Society conferences.
%
% If IEEEtran.cls has not been installed into the LaTeX system files,
% manually specify the path to it like:
% \documentclass[conference]{../sty/IEEEtran}





% Some very useful LaTeX packages include:
% (uncomment the ones you want to load)


% *** MISC UTILITY PACKAGES ***
%
%\usepackage{ifpdf}
% Heiko Oberdiek's ifpdf.sty is very useful if you need conditional
% compilation based on whether the output is pdf or dvi.
% usage:
% \ifpdf
%   % pdf code
% \else
%   % dvi code
% \fi
% The latest version of ifpdf.sty can be obtained from:
% http://www.ctan.org/tex-archive/macros/latex/contrib/oberdiek/
% Also, note that IEEEtran.cls V1.7 and later provides a builtin
% \ifCLASSINFOpdf conditional that works the same way.
% When switching from latex to pdflatex and vice-versa, the compiler may
% have to be run twice to clear warning/error messages.






% *** CITATION PACKAGES ***
%
\usepackage{cite}
% cite.sty was written by Donald Arseneau
% V1.6 and later of IEEEtran pre-defines the format of the cite.sty package
% \cite{} output to follow that of IEEE. Loading the cite package will
% result in citation numbers being automatically sorted and properly
% "compressed/ranged". e.g., [1], [9], [2], [7], [5], [6] without using
% cite.sty will become [1], [2], [5]--[7], [9] using cite.sty. cite.sty's
% \cite will automatically add leading space, if needed. Use cite.sty's
% noadjust option (cite.sty V3.8 and later) if you want to turn this off.
% cite.sty is already installed on most LaTeX systems. Be sure and use
% version 4.0 (2003-05-27) and later if using hyperref.sty. cite.sty does
% not currently provide for hyperlinked citations.
% The latest version can be obtained at:
% http://www.ctan.org/tex-archive/macros/latex/contrib/cite/
% The documentation is contained in the cite.sty file itself.


\usepackage[utf8]{inputenc}


% *** GRAPHICS RELATED PACKAGES ***
%
\ifCLASSINFOpdf
  \usepackage[pdftex]{graphicx}
  % declare the path(s) where your graphic files are
  % \graphicspath{{../pdf/}{../jpeg/}}
  % and their extensions so you won't have to specify these with
  % every instance of \includegraphics
  % \DeclareGraphicsExtensions{.pdf,.jpeg,.png}
\else
  % or other class option (dvipsone, dvipdf, if not using dvips). graphicx
  % will default to the driver specified in the system graphics.cfg if no
  % driver is specified.
  % \usepackage[dvips]{graphicx}
  % declare the path(s) where your graphic files are
  % \graphicspath{{../eps/}}
  % and their extensions so you won't have to specify these with
  % every instance of \includegraphics
  % \DeclareGraphicsExtensions{.eps}
\fi
% graphicx was written by David Carlisle and Sebastian Rahtz. It is
% required if you want graphics, photos, etc. graphicx.sty is already
% installed on most LaTeX systems. The latest version and documentation can
% be obtained at: 
% http://www.ctan.org/tex-archive/macros/latex/required/graphics/
% Another good source of documentation is "Using Imported Graphics in
% LaTeX2e" by Keith Reckdahl which can be found as epslatex.ps or
% epslatex.pdf at: http://www.ctan.org/tex-archive/info/
%
% latex, and pdflatex in dvi mode, support graphics in encapsulated
% postscript (.eps) format. pdflatex in pdf mode supports graphics
% in .pdf, .jpeg, .png and .mps (metapost) formats. Users should ensure
% that all non-photo figures use a vector format (.eps, .pdf, .mps) and
% not a bitmapped formats (.jpeg, .png). IEEE frowns on bitmapped formats
% which can result in "jaggedy"/blurry rendering of lines and letters as
% well as large increases in file sizes.
%
% You can find documentation about the pdfTeX application at:
% http://www.tug.org/applications/pdftex





% *** MATH PACKAGES ***
%
%\usepackage[cmex10]{amsmath}
% A popular package from the American Mathematical Society that provides
% many useful and powerful commands for dealing with mathematics. If using
% it, be sure to load this package with the cmex10 option to ensure that
% only type 1 fonts will utilized at all point sizes. Without this option,
% it is possible that some math symbols, particularly those within
% footnotes, will be rendered in bitmap form which will result in a
% document that can not be IEEE Xplore compliant!
%
% Also, note that the amsmath package sets \interdisplaylinepenalty to 10000
% thus preventing page breaks from occurring within multiline equations. Use:
%\interdisplaylinepenalty=2500
% after loading amsmath to restore such page breaks as IEEEtran.cls normally
% does. amsmath.sty is already installed on most LaTeX systems. The latest
% version and documentation can be obtained at:
% http://www.ctan.org/tex-archive/macros/latex/required/amslatex/math/





% *** SPECIALIZED LIST PACKAGES ***
%
%\usepackage{algorithmic}
% algorithmic.sty was written by Peter Williams and Rogerio Brito.
% This package provides an algorithmic environment fo describing algorithms.
% You can use the algorithmic environment in-text or within a figure
% environment to provide for a floating algorithm. Do NOT use the algorithm
% floating environment provided by algorithm.sty (by the same authors) or
% algorithm2e.sty (by Christophe Fiorio) as IEEE does not use dedicated
% algorithm float types and packages that provide these will not provide
% correct IEEE style captions. The latest version and documentation of
% algorithmic.sty can be obtained at:
% http://www.ctan.org/tex-archive/macros/latex/contrib/algorithms/
% There is also a support site at:
% http://algorithms.berlios.de/index.html
% Also of interest may be the (relatively newer and more customizable)
% algorithmicx.sty package by Szasz Janos:
% http://www.ctan.org/tex-archive/macros/latex/contrib/algorithmicx/




% *** ALIGNMENT PACKAGES ***
%
%\usepackage{array}
% Frank Mittelbach's and David Carlisle's array.sty patches and improves
% the standard LaTeX2e array and tabular environments to provide better
% appearance and additional user controls. As the default LaTeX2e table
% generation code is lacking to the point of almost being broken with
% respect to the quality of the end results, all users are strongly
% advised to use an enhanced (at the very least that provided by array.sty)
% set of table tools. array.sty is already installed on most systems. The
% latest version and documentation can be obtained at:
% http://www.ctan.org/tex-archive/macros/latex/required/tools/


%\usepackage{mdwmath}
%\usepackage{mdwtab}
% Also highly recommended is Mark Wooding's extremely powerful MDW tools,
% especially mdwmath.sty and mdwtab.sty which are used to format equations
% and tables, respectively. The MDWtools set is already installed on most
% LaTeX systems. The lastest version and documentation is available at:
% http://www.ctan.org/tex-archive/macros/latex/contrib/mdwtools/


% IEEEtran contains the IEEEeqnarray family of commands that can be used to
% generate multiline equations as well as matrices, tables, etc., of high
% quality.


%\usepackage{eqparbox}
% Also of notable interest is Scott Pakin's eqparbox package for creating
% (automatically sized) equal width boxes - aka "natural width parboxes".
% Available at:
% http://www.ctan.org/tex-archive/macros/latex/contrib/eqparbox/





% *** SUBFIGURE PACKAGES ***
%\usepackage[tight,footnotesize]{subfigure}
% subfigure.sty was written by Steven Douglas Cochran. This package makes it
% easy to put subfigures in your figures. e.g., "Figure 1a and 1b". For IEEE
% work, it is a good idea to load it with the tight package option to reduce
% the amount of white space around the subfigures. subfigure.sty is already
% installed on most LaTeX systems. The latest version and documentation can
% be obtained at:
% http://www.ctan.org/tex-archive/obsolete/macros/latex/contrib/subfigure/
% subfigure.sty has been superceeded by subfig.sty.



%\usepackage[caption=false]{caption}
%\usepackage[font=footnotesize]{subfig}
% subfig.sty, also written by Steven Douglas Cochran, is the modern
% replacement for subfigure.sty. However, subfig.sty requires and
% automatically loads Axel Sommerfeldt's caption.sty which will override
% IEEEtran.cls handling of captions and this will result in nonIEEE style
% figure/table captions. To prevent this problem, be sure and preload
% caption.sty with its "caption=false" package option. This is will preserve
% IEEEtran.cls handing of captions. Version 1.3 (2005/06/28) and later 
% (recommended due to many improvements over 1.2) of subfig.sty supports
% the caption=false option directly:
%\usepackage[caption=false,font=footnotesize]{subfig}
%
% The latest version and documentation can be obtained at:
% http://www.ctan.org/tex-archive/macros/latex/contrib/subfig/
% The latest version and documentation of caption.sty can be obtained at:
% http://www.ctan.org/tex-archive/macros/latex/contrib/caption/




% *** FLOAT PACKAGES ***
%
%\usepackage{fixltx2e}
% fixltx2e, the successor to the earlier fix2col.sty, was written by
% Frank Mittelbach and David Carlisle. This package corrects a few problems
% in the LaTeX2e kernel, the most notable of which is that in current
% LaTeX2e releases, the ordering of single and double column floats is not
% guaranteed to be preserved. Thus, an unpatched LaTeX2e can allow a
% single column figure to be placed prior to an earlier double column
% figure. The latest version and documentation can be found at:
% http://www.ctan.org/tex-archive/macros/latex/base/



%\usepackage{stfloats}
% stfloats.sty was written by Sigitas Tolusis. This package gives LaTeX2e
% the ability to do double column floats at the bottom of the page as well
% as the top. (e.g., "\begin{figure*}[!b]" is not normally possible in
% LaTeX2e). It also provides a command:
%\fnbelowfloat
% to enable the placement of footnotes below bottom floats (the standard
% LaTeX2e kernel puts them above bottom floats). This is an invasive package
% which rewrites many portions of the LaTeX2e float routines. It may not work
% with other packages that modify the LaTeX2e float routines. The latest
% version and documentation can be obtained at:
% http://www.ctan.org/tex-archive/macros/latex/contrib/sttools/
% Documentation is contained in the stfloats.sty comments as well as in the
% presfull.pdf file. Do not use the stfloats baselinefloat ability as IEEE
% does not allow \baselineskip to stretch. Authors submitting work to the
% IEEE should note that IEEE rarely uses double column equations and
% that authors should try to avoid such use. Do not be tempted to use the
% cuted.sty or midfloat.sty packages (also by Sigitas Tolusis) as IEEE does
% not format its papers in such ways.





% *** PDF, URL AND HYPERLINK PACKAGES ***
%
%\usepackage{url}
% url.sty was written by Donald Arseneau. It provides better support for
% handling and breaking URLs. url.sty is already installed on most LaTeX
% systems. The latest version can be obtained at:
% http://www.ctan.org/tex-archive/macros/latex/contrib/misc/
% Read the url.sty source comments for usage information. Basically,
% \url{my_url_here}.





% *** Do not adjust lengths that control margins, column widths, etc. ***
% *** Do not use packages that alter fonts (such as pslatex).         ***
% There should be no need to do such things with IEEEtran.cls V1.6 and later.
% (Unless specifically asked to do so by the journal or conference you plan
% to submit to, of course. )


% correct bad hyphenation here
\hyphenation{op-tical net-works semi-conduc-tor}


\begin{document}
%
% paper title
% can use linebreaks \\ within to get better formatting as desired
\title{Proposal for True Random Generation on the APQ8060 Dragonboard}

% author names and affiliations
% use a multiple column layout for up to three different
% affiliations
\author{
\IEEEauthorblockN{Kevin Burns}
\IEEEauthorblockA{Department of Electrical and\\Computer Engineering\\
Virginia Tech\\
Blacksburg, Virginia\\
Email: kevinpb@vt.edu}
\and
\IEEEauthorblockN{Rob Lyerly}
\IEEEauthorblockA{Department of Electrical and\\Computer Engineering\\
Virginia Tech\\
Blacksburg, Virginia\\
Email: rlyerly@vt.edu}
\and
\IEEEauthorblockN{Reese Moore}
\IEEEauthorblockA{Department of Electrical and\\Computer Engineering\\
Virginia Tech\\
Blacksburg, Virginia\\
Email: ram@vt.edu}
\and
\IEEEauthorblockN{Philip Kobezak}
\IEEEauthorblockA{Department of Electrical and\\Computer Engineering\\
Virginia Tech\\
Blacksburg, Virginia\\
Email: pkobezak@vt.edu}
}

% conference papers do not typically use \thanks and this command
% is locked out in conference mode. If really needed, such as for
% the acknowledgment of grants, issue a \IEEEoverridecommandlockouts
% after \documentclass

% for over three affiliations, or if they all won't fit within the width
% of the page, use this alternative format:
% 
%\author{\IEEEauthorblockN{Michael Shell\IEEEauthorrefmark{1},
%Homer Simpson\IEEEauthorrefmark{2},
%James Kirk\IEEEauthorrefmark{3}, 
%Montgomery Scott\IEEEauthorrefmark{3} and
%Eldon Tyrell\IEEEauthorrefmark{4}}
%\IEEEauthorblockA{\IEEEauthorrefmark{1}School of Electrical and Computer Engineering\\
%Georgia Institute of Technology,
%Atlanta, Georgia 30332--0250\\ Email: see http://www.michaelshell.org/contact.html}
%\IEEEauthorblockA{\IEEEauthorrefmark{2}Twentieth Century Fox, Springfield, USA\\
%Email: homer@thesimpsons.com}
%\IEEEauthorblockA{\IEEEauthorrefmark{3}Starfleet Academy, San Francisco, California 96678-2391\\
%Telephone: (800) 555--1212, Fax: (888) 555--1212}
%\IEEEauthorblockA{\IEEEauthorrefmark{4}Tyrell Inc., 123 Replicant Street, Los Angeles, California 90210--4321}}




% use for special paper notices
%\IEEEspecialpapernotice{(Invited Paper)}




% make the title area
\maketitle


\begin{abstract}
%\boldmath
    With the mobile platform becoming more ubiquitous in modern society there is a greater need for the ability to run cryptographic algorithms on 
    low power architectures. Most of these algorithms require a source of randomness. For the Android version of the Linux kernel this is provided
    by user input and timing data. It have been observed that this methodology does not provide the user with enough randomness and has caused either
    system performance issues or none-unique cryptographic keys. Therefore there is a high demand for a quick and power efficient random number
    generator for mobile devices.
\end{abstract}

% IEEEtran.cls defaults to using nonbold math in the Abstract.
% This preserves the distinction between vectors and scalars. However,
% if the conference you are submitting to favors bold math in the abstract,
% then you can use LaTeX's standard command \boldmath at the very start
% of the abstract to achieve this. Many IEEE journals/conferences frown on
% math in the abstract anyway.

% no keywords

% For peer review papers, you can put extra information on the cover
% page as needed:
% \ifCLASSOPTIONpeerreview
% \begin{center} \bfseries EDICS Category: 3-BBND \end{center}
% \fi
%
% For peerreview papers, this IEEEtran command inserts a page break and
% creates the second title. It will be ignored for other modes.
\IEEEpeerreviewmaketitle



\section{Introduction}


\subsection{Related Work and Background}

%Seeder
In the Linux kernel, including the Android kernel, there is an entropy pool. This pool is populated via environmental noise from device drivers and other 
sources into an entropy pool. These sources are treated as true random number generation and therefore create entropy suitable for key generation and other
cryphotgraphy. Access to this pool is granted through the random device located in /dev/random \cite{random_man} and through the urandom device, 
located in /dev/urandom. These two interfaces read straight from the pool until a limit of 128 bits of remaining entropy is reached. If the pool is drained to below this bit limit, and if the application reading from the pool is using the random interface, then it will spin until more entropy is added to the pool. If the application is using the urandom interface, which provides ``unlimited'' entropy. The urandom device does not block, but forgoes true randomness and uses and feeds psuedo random data into the entropy pool. It is stated that this method is theoretically vulnerable to cryptographic attacks. 

In the Google Play Store there exists an application entitled Seeder. The developers of seeder are trying to speed up the Android experience by removing
the need to ever spin on the random device. Their solution to provide entropy even when the random device's 128 bit limit is reached by to piping the 
false entropy from the urandom device back into the entropy pool. The developers of Seeder claim this brings the user noticeable speedups in day to 
day use of the Android Operating System. As this application has been downloaded over ten thousand times in the past 30 days (as of the date this paper was written)
, it is clear that this is an area of interest to the Android user community.

%The Sources of Randomness in Mobile Devices
The task of finding true randomness on mobile platforms has been targeted before. In \cite{Krhovjak}, the two main sources analyzed were the 
camera and microphone. The authors of this paper were unable to generate definite rates for random number generation but were able to verify the validity.
There were several advantages and disadvantages presented in this paper in regards to the camera as a random number generator. The advantages included the
ability to produce considerable entropy even at $5\,^{\circ}\mathrm{C}$ and when the camera guard was closed. However there were also some vulnerabilities, 
an attacker could use a halogen light bulb to saturate the CCD to produce only 255 values for R, G, and B values of each pixel. This would highly bias the randomness produced by the camera. 
The microphone's noise as a source of randomness was found to be 3 bits of entropy per sample given my the Shannon entropy formula and 0.5 using the min-entropy 
formula. The only attempt to skew the data was using two sets of recordings one having music played into it and one with ambient background noises. They did
find correlations between randomness generated from the music recording, but stated that this could be corrected with the proper post-processing. The authors
gave no mention to the power draw of each method. It was also mentioned that manufacturers of both microphones cameras are continuously trying to decrease the noise
via both hardware and proprietary (hidden) software. The tests mentioned in \cite{Krhovjak} were performed in 2007 on dated hardware.

%Accelerometers and Randomness: Perfect Together
Accelerometers have been studied as a potential source of entropy in low-power and embedded devices.  J. Voris et al. analyzed the use of accelerometers in an
Intel WISP RFID tag \cite{voris}.  They conclude that even when stationary, accelerometers are sensitive enough to provide a good rate of entropy; they even
test the device when at rest on industrial dampening material.  They show that accelerometers are highly resistant to bias, arguing that the best attack
vector to reduce randomness is to leave the accelerometer stationary (i.e. movement only serves to increase the randomness rather than degrade it).  While it is
conceivable that an attacker may be able to saturate the device's output by applying an extraordinary force (say, in a centrifuge), an accelerometer lacks the
potential for bias that plagues other devices on the phone.  Although gyroscopes have not been studied as a means for random number generation, they are a
similar device in principal to the accelerometer, striving to achieve a high level of sensitivity.  In addition, these two devices have become ubiquitous on
modern phones and are extremely cheap to include in any system.  Because of the low potential for bias and the widespread use of these devices, they are an
attractive means for adding entropy to a random number generator.

\subsection{Problem Statement}
Due to the increasing number of mobile platforms and their use of cryptographic algorithms and random number generators, there is demand for quick 
and power efficient random number generators in the mobile environment. Specifically on our target hardware, Qualcomm's APQ8060 Dragonboard, we have 
access to the random and urandom devices given via the Linux kernel. The random device is the true random number generation option in the Linux 
environment, it is therefore a safe baseline for comparison as it is fairly consistent between Android devices. There are issues with this current solution.
Firstly, a performance decrease can be witnessed when the entropy pool, used by random, is depleted to 128 bits. This is may be more prevalent in Android's
Event Loop software model, as there isn't a clear multi-tasking structure. Another flaw comes to the scene when programmers and users try to circumvent
this performance drain, by either filling the entropy pool with a none random bit-stream or by using the urandom device where the random device is more secure.
We witness the effects of these issues in sources like \cite{ron_was_wrong}, which found that 4\% of 6.6 million certificates and PGP keys use the same RSA 
modulus and therefore possibly have the exact same private key as someone else. The authors of \cite{chip_and_skim}, exposed poor random number generators
in EMV machines. Companies, such as Qualcomm, provide hardware random number generators in their System on a Chip (SoC) architectures to combat these drawbacks. 
Therefore there is a clear need for a low power solution that is applicable to existing and future mobile devices on the market.


\section{Proposed Solution}
\begin{figure}[t]
	\includegraphics[width=\columnwidth]{proposed_solution}
	\caption{Block diagram of the proposed solution.  Our daemon will poll from the accelerometer and gyroscope (when requested), perform post-processing and feed random bits into the Linux entropy pool via system calls.}
	\label{proposed_solution_block}
\end{figure}

Although mobile phones have many different devices on-chip that could be potentially used as sources of randomness for a random number generator,
most of these devices have the potential for strong bias.  Additionally, because of the operating environment polling these devices with high
frequency can cause serious system performance and energy issues.  Therefore it is necessary to not only use devices that have high quality randomness,
but to use those devices in an intelligent manner so as to not overwhelm the computational and energy resources of the device.

We propose the implementation of a daemon process that polls data from the accelerometer and gyroscope (when the entropy pool is running low), processes
input data to eliminate sources of bias, and feeds the data into the entropy pool.  Figure \ref{proposed_solution_block} shows the architecture for
our daemon.  Android built on Linux exposes standard interfaces to the hardware sensors for data retrieval; additionally, Linux provides system calls
to feed random data into the entropy pool.  Our daemon will interact with both of these interfaces.

There are many types of post processing that can be implemented to increase the uniformity of the random numbers generated.  Hash algorithms like SHA-3
\cite{keccak} can be used to eliminate bias.  Similarly, Barak et al. propose a technique whereby random numbers are pulled from several independent
devices \cite{independent_devices}.  Two of these numbers are multiplied together and added to a third number to strengthen the randomness.  Barak et
al. propose another technique where random numbers are multiplied by a Toeplitz matrix to approach a uniform distribution \cite{true_rng}.  These
techniques will be implemented and compared for their effectiveness and computational complexity.

Random numbers generated by our daemon will be compared to the stock Android/Linux random number generator (which includes an on-board hardware
RNG that may not be present on all mobile devices).  These two sources will be compared on three dimensions: quality of random numbers generated,
energy draw from using the random number generator, and computational resources used by random number generation.  The DIEHARD \cite{diehard}/NIST 
\cite{nist} benchmarks will be used to evaluate the quality of the random numbers generated; synthetic benchmarks will be written to evaluate the energy
draw and computational complexity of the random number generators.

%% Talk about the power usage
% Power Measurements

Efficient generation of random numbers, with high entropy, is ensential 
for mobile devices.  Power must always be used carefully as to not drain 
the battery too quickly.  Any peripheral device, such as an 
accelerometer, will use more power when it is collecting data.  The 
APQ8060 Dragonboard uses a Kionix KXSD9-1026 accelerometer which draws 
typically 0.3 \mu Amps while in standby \cite{Kionix}.  However, full 
operating power draws more than 200 \mu Amps.  Running the accelerometer 
for any length of time, at a high sampling rate, may draw more power 
than other sources of entropy.  For this task of the project, the 
implemented daemon will be run while sampling the current draw of the 
whole system.  The result will show if the proposed solution has a 
significant increase in power usage over the stock RNG.


\subsection{Expectations}
Our expectations are to successfully produce a bit stream generated using the accelerometer and gyroscope to meet the requirements tested in NIST and DIEHARD.
Most notably we would like to achieve a normal distribution of random numbers that are suitable for adding to the Linux kernel's entropy pool. We also expect
to achieve a fast enough rate to run only when the entropy is at critical levels, thus saving power.

%% Expectations

From the findings of \cite{voris}, accelerometers appear to be highly random. A
high rate of sampling shows that they are fairly fast and should be able to
provide a high rate of random bits per second when they are needed. The
specification sheet for the accelerometer found on the Dragonboard shows that it
is also efficient for power consumption.

The apparent randomness of the accelerometer and gyroscope, accompanied with the
relatively low power nature of the components as well as their ubiquity on smart
phones, tends to suggest that they would be an excellent source of entropy for
mobile platforms. A solution that pulls from these sensors should be easily
deployed to several different platforms, including those that do not have the
benefit of an on board hardware random number generator. As quality random
numbers are necessary for most of the common cryptographic operations, such a
system will be able to bring faster secure communication to more devices. 


\section{Plan for Development}
%% Tasks



\subsection{Risks and Unknowns}
%% Risks and Unknowns

% Talk about:
%  - Is the accelerometer really entropic?
%  - Will this be worthwhile vs. the hw-rng?
%  - Is this data being incorporated somehow already?
%  - Unknown whether post processing is ``good''.

For any project, it is important to understand that there are risks and unknowns
that will need to be understood and mitigated during the course of the project.
By assessing what may become a problem during the process of the project, those
problems can be consciously avoided or dealt with early on.

One unknown that is central to this project is the entropy of the accelerometer.
Based on the findings in \cite{voris}, it appears that accelerometers are good
sources of entropy, it is possible that the accelerometer on the Dragonboard
violates some assumption or that technology has changed such that modern
accelerometers are no longer good sources of entropy. This can be tested using
several randomness tests, such as the DIEHARD and NIST entropy tests, which will
be essential to do before moving on with any sort of analysis.  These tests will
show if the sensor data appears to have good entropy or not.  In the case that
it does not, it would invalidate the remaining work for the project. 

A potential problem is the existence of the hardware random number generator
that is already on the board. Depending on how good that hardware is, it could
potentially dwarf any results that this project generates. Even so, the results
should be valuable, both as a look at the prospect of using accelerometers in
future designs for hardware random number generators as well as for providing
better entropy to existing platforms which do not have hardware random number
generators. 

During the testing phase of the project, members of the project team will
perform an experiment to show what sort of data is experienced over the course
of a couple of days. As it is not reasonable to carry the Dragonboard over the
course of a daily routine, some of the data will need to be collected on
personal cell phones. It could potentially be a problem if no available cell
phone has a similar accelerometer or gyroscope as this would introduce an
additional variable into the analysis. Research and experimentation will have to
be done to discover if, and show, the sensors in the selected cell phones will
be a reasonable stand in for the Dragonboard. If nothing matches closely enough,
a different method of collecting longer term real world data will have to be
devised.

If the data from these sensors is already being incorporated into the entropy
pool, it would make this work largely a repetition of that previous work. Upon
first inspection this does not appear to be the case, though further research is
needed to ensure that this is indeed novel work.


\section{Deliverables}
% Deliverables

The deliverables of this project will include a detailed report, the 
daemon source, post processing source, and any scripts used in 
development. Additionally, the results will be presented, summarizing 
the findings in the report.  

The results of this work should show that the potential exists to use 
common peripherals on mobile devices to help fill the entropy pool of 
the RNG. This will be beneficial to developers who want to ensure a 
consistent stream of high-quality random numbers. Even though the 
end-user of these devices may not be aware of the RNG, they will notice 
slow downs in encrypted application sessions. The proposed solution 
should reduce the chances of end-users experiencing slow downs when 
concurrent encrypted sessions require random numbers quickly.




% An example of a floating figure using the graphicx package.
% Note that \label must occur AFTER (or within) \caption.
% For figures, \caption should occur after the \includegraphics.
% Note that IEEEtran v1.7 and later has special internal code that
% is designed to preserve the operation of \label within \caption
% even when the captionsoff option is in effect. However, because
% of issues like this, it may be the safest practice to put all your
% \label just after \caption rather than within \caption{}.
%
% Reminder: the "draftcls" or "draftclsnofoot", not "draft", class
% option should be used if it is desired that the figures are to be
% displayed while in draft mode.
%
%\begin{figure}[!t]
%\centering
%\includegraphics[width=2.5in]{myfigure}
% where an .eps filename suffix will be assumed under latex, 
% and a .pdf suffix will be assumed for pdflatex; or what has been declared
% via \DeclareGraphicsExtensions.
%\caption{Simulation Results}
%\label{fig_sim}
%\end{figure}

% Note that IEEE typically puts floats only at the top, even when this
% results in a large percentage of a column being occupied by floats.


% An example of a double column floating figure using two subfigures.
% (The subfig.sty package must be loaded for this to work.)
% The subfigure \label commands are set within each subfloat command, the
% \label for the overall figure must come after \caption.
% \hfil must be used as a separator to get equal spacing.
% The subfigure.sty package works much the same way, except \subfigure is
% used instead of \subfloat.
%
%\begin{figure*}[!t]
%\centerline{\subfloat[Case I]\includegraphics[width=2.5in]{subfigcase1}%
%\label{fig_first_case}}
%\hfil
%\subfloat[Case II]{\includegraphics[width=2.5in]{subfigcase2}%
%\label{fig_second_case}}}
%\caption{Simulation results}
%\label{fig_sim}
%\end{figure*}
%
% Note that often IEEE papers with subfigures do not employ subfigure
% captions (using the optional argument to \subfloat), but instead will
% reference/describe all of them (a), (b), etc., within the main caption.


% An example of a floating table. Note that, for IEEE style tables, the 
% \caption command should come BEFORE the table. Table text will default to
% \footnotesize as IEEE normally uses this smaller font for tables.
% The \label must come after \caption as always.
%
%\begin{table}[!t]
%% increase table row spacing, adjust to taste
%\renewcommand{\arraystretch}{1.3}
% if using array.sty, it might be a good idea to tweak the value of
% \extrarowheight as needed to properly center the text within the cells
%\caption{An Example of a Table}
%\label{table_example}
%\centering
%% Some packages, such as MDW tools, offer better commands for making tables
%% than the plain LaTeX2e tabular which is used here.
%\begin{tabular}{|c||c|}
%\hline
%One & Two\\
%\hline
%Three & Four\\
%\hline
%\end{tabular}
%\end{table}


% Note that IEEE does not put floats in the very first column - or typically
% anywhere on the first page for that matter. Also, in-text middle ("here")
% positioning is not used. Most IEEE journals/conferences use top floats
% exclusively. Note that, LaTeX2e, unlike IEEE journals/conferences, places
% footnotes above bottom floats. This can be corrected via the \fnbelowfloat
% command of the stfloats package.



%\section{Conclusion}
%The conclusion goes here.




% conference papers do not normally have an appendix


% use section* for acknowledgement
%\section*{Acknowledgment}


%The authors would like to thank...

% trigger a \newpage just before the given reference
% number - used to balance the columns on the last page
% adjust value as needed - may need to be readjusted if
% the document is modified later
%\IEEEtriggeratref{8}
% The "triggered" command can be changed if desired:
%\IEEEtriggercmd{\enlargethispage{-5in}}

% references section

% can use a bibliography generated by BibTeX as a .bbl file
% BibTeX documentation can be easily obtained at:
% http://www.ctan.org/tex-archive/biblio/bibtex/contrib/doc/
% The IEEEtran BibTeX style support page is at:
% http://www.michaelshell.org/tex/ieeetran/bibtex/
%\bibliographystyle{IEEEtran}
% argument is your BibTeX string definitions and bibliography database(s)
%\bibliography{IEEEabrv,../bib/paper}
%
% <OR> manually copy in the resultant .bbl file
% set second argument of \begin to the number of references
% (used to reserve space for the reference number labels box)


\bibliographystyle{plain}
\bibliography{phase2}


% that's all folks
\end{document}


